\documentclass[12pt,a4paper]{article}

% Packages
\usepackage[utf8]{inputenc}
\usepackage[margin=1in]{geometry}
\usepackage{graphicx}
\usepackage{float}
\usepackage{booktabs}
\usepackage{longtable}
\usepackage{array}
\usepackage{hyperref}
\usepackage{xcolor}
\usepackage{listings}
\usepackage{amsmath}
\usepackage{amssymb}
\usepackage{caption}
\usepackage{subcaption}
\usepackage{fancyhdr}
\usepackage{titlesec}
\usepackage{enumitem}
\usepackage{multirow}
\usepackage[skip=10pt]{parskip}

% Hyperref setup
\hypersetup{
    colorlinks=true,
    linkcolor=blue,
    filecolor=magenta,
    urlcolor=cyan,
    citecolor=green,
    pdftitle={Gem5 MergeSort Cache Analysis},
    pdfauthor={Sumit Kumar, Aviral Singh},
    pdfsubject={High Performance in Computer Architecture},
}

% Code listing setup
\lstset{
    basicstyle=\ttfamily\small,
    breaklines=true,
    frame=single,
    numbers=left,
    numberstyle=\tiny,
    commentstyle=\color{gray},
    keywordstyle=\color{blue},
    stringstyle=\color{red}
}

% Header and footer
\pagestyle{fancy}
\fancyhf{}
\rhead{High Performance in Computer Architecture}
\lhead{Problem 2 - Assignment 1}
\cfoot{Page \thepage}

% Title formatting
\titleformat{\section}{\Large\bfseries}{\thesection}{1em}{}
\titleformat{\subsection}{\large\bfseries}{\thesubsection}{1em}{}
\titleformat{\subsubsection}{\normalsize\bfseries}{\thesubsubsection}{1em}{}

\begin{document}

% ============================================================================
% TITLE PAGE
% ============================================================================
\begin{titlepage}
    \centering
    \vspace*{1.5cm}

    {\Huge\bfseries Gem5 MergeSort Cache Analysis\par}
    \vspace{0.5cm}
    {\Large Cache Optimization for Memory-Intensive Workloads\par}
    \vspace{1.5cm}

    {\LARGE Problem 2: Assignment 1\par}
    \vspace{2.5cm}

    {\Large\bfseries Submitted by:\par}
    \vspace{0.5cm}
    {\Large Sumit Kumar\par}
    {\large Roll No: 22CS30056\par}
    \vspace{0.5cm}
    {\Large Aviral Singh\par}
    {\large Roll No: 22CS30015\par}
    \vspace{3cm}

    {\large\bfseries Course:\par}
    {\Large High Performance in Computer Architecture\par}
    \vspace{0.5cm}
    {\large Department of Computer Science and Engineering\par}

    \vfill
    {\large \today\par}
\end{titlepage}

% ============================================================================
% ABSTRACT
% ============================================================================
\newpage
\thispagestyle{empty}
\begin{abstract}
This report presents a comprehensive analysis of cache memory optimization for merge sort algorithms using the gem5 simulator with RISC-V ISA. We compare two merge sort implementations: a simple in-memory version operating on a 10MB dataset and a chunked version processing 2MB chunks with 1MB streaming buffers. Through extensive simulations totaling 164 configurations, we analyze the impact of L1 and L2 cache size and associativity on performance metrics including IPC, miss rates, and execution time.

Our findings demonstrate that the chunked algorithm consistently outperforms the simple version by 2\% across all cache configurations, achieving a best IPC of 0.2949 with 128KiB L1 and 1024KiB L2 caches (both 4-way associative). The study reveals that L2 cache capacity is the dominant performance factor, while L1 associativity has minimal impact due to the sequential access patterns inherent in merge sort algorithms. Surprisingly, higher L2 associativity (16-way) slightly degrades performance compared to 4-way, likely due to increased tag comparison overhead without benefit for streaming access patterns.

Key contributions of this work include: (1) systematic exploration of 81 cache configurations across two algorithms, (2) identification of optimal cache parameters for sorting workloads, (3) quantification of algorithm-cache interaction effects, and (4) validation that cache-aware algorithm design provides consistent benefits across diverse hardware configurations.
\end{abstract}

% ============================================================================
% TABLE OF CONTENTS
% ============================================================================
\newpage
\tableofcontents

% ============================================================================
% SECTION 1: INTRODUCTION
% ============================================================================
\newpage
\section{Introduction}

This assignment explores cache memory hierarchy optimization for memory-intensive sorting workloads using the gem5 architectural simulator. Cache memory plays a critical role in bridging the performance gap between fast processors and relatively slow main memory. 

\subsection{Assignment Overview}

The assignment consists of two main components:

\textbf{Part 1: Baseline Comparison}
\begin{itemize}
    \item Compare two merge sort variants using default cache configurations
    \item Analyze cache behavior and identify bottlenecks
    \item Understand why chunked sorting shows better cache locality
\end{itemize}

\textbf{Part 2: Cache Optimization Sweep}
\begin{itemize}
    \item Test 81 different cache configurations for each algorithm
    \item Identify optimal L1 and L2 cache parameters
    \item Analyze performance trade-offs between size and associativity
    \item Generate comprehensive visualizations and analysis
\end{itemize}

\subsection{Merge Sort Algorithms}

\subsubsection{Simple Merge Sort}

The simple merge sort operates on the entire 10MB dataset (2,621,440 integers) in memory using a standard recursive divide-and-conquer approach. It reads from the \texttt{random\_numbers.bin} file, sorts the data in-place, and writes the sorted output.

\textbf{Characteristics:}
\begin{itemize}
    \item Working set: 10MB (entire dataset)
    \item Memory access pattern: Recursive, in-place merging
    \item Cache behavior: Frequent evictions due to large working set
    \item Binary size: 497KB (mergesort\_s)
\end{itemize}

\subsubsection{Chunked Merge Sort}

The chunked merge sort employs the approach:

\textbf{Phase 1 - Chunking:} Divides 10MB dataset into 5 chunks of 2MB each

\textbf{Phase 2 - Local Sorting:} Sorts each 2MB chunk independently using standard merge sort

\textbf{Phase 3 - Streaming Merge:} Performs 5-way parallel merge using 1MB input streams and 1MB output buffer

\textbf{Characteristics:}
\begin{itemize}
    \item Working set: 2MB per chunk during sorting, 1MB streams during merge
    \item Memory access pattern: Sequential streaming with better locality
    \item Cache behavior: More predictable prefetching opportunities
    \item Binary size: 502KB (mergesort\_c)
\end{itemize}

\subsection{Simulation Environment}

\textbf{Gem5 Simulator Configuration:}
\begin{itemize}
    \item ISA: RISC-V 64-bit 
    \item CPU Model: TimingSimpleCPU
    \item Clock Frequency: 1 GHz
    \item Memory: 2GB DDR3
    \item Cache Line Size: 64 bytes
\end{itemize}

\textbf{Default Cache Configuration (Part 1):}
\begin{table}[H]
\centering
\begin{tabular}{lcccc}
\toprule
\textbf{Cache Level} & \textbf{Size} & \textbf{Associativity} \\
\midrule
L1 Instruction & 32 KiB & 8-way  \\
L1 Data & 64 KiB & 8-way & \\
L2 Unified & 512 KiB & 16-way  \\
\bottomrule
\end{tabular}
\caption{Default cache configuration used in Part 1 baseline analysis}
\label{tab:default_config}
\end{table}

\subsection{Simulation Scale}

\textbf{Total Simulations Conducted:} 164
\begin{itemize}
    \item Part 1 Baseline: 2 simulations
    \item Part 2 Cache Sweep: 162 simulations (81 configurations $\times$ 2 algorithms)
    \item Total Execution Time: Approximately 11 hours
\end{itemize}

% ============================================================================
% SECTION 2: PART 1 - BASELINE COMPARISON
% ============================================================================
\newpage
\section{Part 1: Baseline Comparison}

Part 1 establishes baseline performance metrics for both merge sort algorithms using the default cache configuration. 

\subsection{Baseline Statistics Comparison}

Table \ref{tab:baseline_stats} presents a comprehensive comparison of performance metrics between the simple and chunked merge sort implementations under default cache settings.

\begin{table}[H]
\centering
\small
\begin{tabular}{p{4cm}rrrc}
\toprule
\textbf{Metric} & \textbf{Simple} & \textbf{Chunked} & \textbf{Difference} & \textbf{Better Performance} \\
\midrule
\multicolumn{5}{l}{\textbf{Execution Metrics}} \\
sim\_ticks & 16,632,464,767,000 & 17,264,476,070,000 & +3.8\% & Simple \\
sim\_insts & 4,788,525,072 & 5,067,505,484 & +5.8\% & Simple \\
\textbf{IPC} & \textbf{0.288} & \textbf{0.294} & \textbf{+2.0\%} & \textbf{Chunked} \\
\midrule
\multicolumn{5}{l}{\textbf{L1D Cache (Data)}} \\
Demand Accesses & 2,271,121,566 & 2,331,978,469 & +2.7\% & Simple \\
Demand Hits & 2,265,238,366 & 2,327,230,897 & +2.7\% & Chunked \\
Demand Misses & 5,883,200 & 4,747,572 & -19.3\% & Chunked \\
\textbf{Miss Rate} & \textbf{0.259\%} & \textbf{0.204\%} & \textbf{-21.4\%} & \textbf{Chunked} \\
\midrule
\multicolumn{5}{l}{\textbf{L1I Cache (Instruction)}} \\
Demand Accesses & 5,772,498,057 & 6,112,641,646 & +5.9\% & Simple \\
Demand Misses & 453 & 514 & +13.5\% & Simple \\
Miss Rate & $\sim$0\% & $\sim$0\% & Negligible & Both \\
\midrule
\multicolumn{5}{l}{\textbf{L2 Cache (Unified)}} \\
Demand Accesses & 5,883,658 & 4,748,089 & -19.3\% & Chunked \\
Demand Hits & 1,979,245 & 2,040,366 & +3.1\% & Chunked \\
Demand Misses & 3,904,413 & 2,707,723 & -30.6\% & Chunked \\
\textbf{Miss Rate} & \textbf{66.4\%} & \textbf{57.0\%} & \textbf{-14.1\%} & \textbf{Chunked} \\
\bottomrule
\end{tabular}
\caption{Baseline performance comparison with default cache configuration}
\label{tab:baseline_stats}
\end{table}

\subsection{Key Observations}

\subsubsection{1. Chunked Algorithm Shows Superior Cache Performance}

Despite executing \textbf{5.8\% more instructions} and consuming \textbf{3.8\% more cycles}, the chunked version achieves:

\begin{itemize}
    \item 2.0\% higher IPC (0.294 vs 0.288)
    \item 21.4\% lower L1D miss rate (0.204\% vs 0.259\%)
    \item 30.6\% fewer L2 misses (2.7M vs 3.9M)
    \item 14.1\% lower L2 miss rate (57.0\% vs 66.4\%)
\end{itemize}

This demonstrates that \textbf{cache efficiency can overcome algorithmic overhead}.

\subsubsection{2. L1 Instruction Cache is Near-Perfect}

Both algorithms exhibit L1I hit rate greater than 99.99999\% with only approximately 500 instruction misses across billions of accesses. This indicates highly predictable loop-based instruction patterns and validates our decision in Part 2 to focus optimization on L1D and L2 caches.

\subsubsection{3. L2 Cache is the Primary Bottleneck}

Both algorithms show very high L2 miss rates (57-66\%) while L1D performs well (greater than 99.7\% hit rate). The 512 KiB L2 is insufficient for the 10MB working set, motivating testing of larger L2 configurations in Part 2.

\subsection{Analysis: Why Chunked Sorting Shows Better Cache Locality}

The chunked merge sort demonstrates superior cache performance despite its algorithmic complexity, primarily due to improved temporal and spatial locality.

\subsubsection{Working Set Size Advantage}

The simple merge sort operates on the entire 10MB dataset simultaneously, causing frequent cache evictions as the working set (10MB) far exceeds both L1D (64 KiB) and L2 (512 KiB) capacities. During recursive merging, data accessed early in the algorithm is repeatedly evicted and reloaded, leading to poor cache reuse.

In contrast, the chunked version processes 2MB chunks that still exceed cache capacity but benefit from a critical optimization: the 5-way parallel merge phase uses 1MB streams. Each stream's 1MB buffer (262,144 integers) creates a more cache-friendly access pattern. While individual streams still exceed cache sizes, the algorithm sequentially reads from each stream, maintaining better sequential access patterns that leverage cache line prefetching and minimize thrashing.

\subsubsection{Quantitative Cache Locality Benefits}

The results clearly show chunked sorting's advantages in specific parameters:

\begin{itemize}
    \item \textbf{L1D Miss Rate:} 0.204\% (chunked) vs 0.259\% (simple) -- 21.4\% improvement
    \item \textbf{L2 Miss Rate:} 57.0\% (chunked) vs 66.4\% (simple) -- 14.1\% improvement
    \item \textbf{L2 Demand Misses:} 2.7M (chunked) vs 3.9M (simple) -- 30.6\% reduction
    \item \textbf{IPC:} 0.294 (chunked) vs 0.288 (simple) -- 2.0\% improvement
\end{itemize}

\subsubsection{Hardware Prefetcher Synergy}

The chunked approach's streaming merge pattern creates more predictable memory access, allowing the processor's cache prefetcher to work effectively. Sequential reads from fixed-size buffers enable the hardware to anticipate memory requests and preload cache lines before they are accessed. Although the chunked version executes 5.8\% more instructions due to chunk management overhead, the improved cache hit rates more than compensate, yielding better overall performance as measured by IPC.

% ============================================================================
% SECTION 3: PART 2 - CACHE OPTIMIZATION SWEEP
% ============================================================================
\newpage
\section{Part 2: Cache Optimization Sweep}

Part 2 explores the cache parameter space systematically to identify optimal configurations for both merge sort algorithms.

\subsection{Methodology}

\subsubsection{Parameter Space Design}

We tested 81 unique cache configurations for each algorithm, totaling 162 simulations.
Note: Here the values of L1I and L1D cache are kept the same always.

\begin{table}[H]
\centering
\begin{tabular}{lcc}
\toprule
\textbf{Parameter} & \textbf{Values Tested} & \textbf{Count} \\
\midrule
L1I Size & 32 KiB, 64 KiB, 128 KiB & 3 \\
L1D Size & 32 KiB, 64 KiB, 128 KiB & 3 \\
L1 Associativity & 4-way, 8-way, 16-way & 3 \\
L2 Size & 256 KiB, 512 KiB, 1024 KiB & 3 \\
L2 Associativity & 4-way, 8-way, 16-way & 3 \\
\midrule
\multicolumn{2}{l}{\textbf{Total Configurations}} & \textbf{81} \\
\multicolumn{2}{l}{\textbf{Total Simulations}} & \textbf{162} \\
\bottomrule
\end{tabular}
\caption{Cache parameter space configuration}
\label{tab:param_space}
\end{table}

\subsubsection{Design Decision: L1I and L1D Symmetry}

We configured L1I and L1D caches with identical size and associativity in each configuration. This decision is justified by:

\begin{enumerate}
    \item \textbf{Baseline Analysis:} L1I shows near-perfect hit rates (greater than 99.99999\%), indicating it is not a bottleneck
    \item \textbf{Architectural Realism:} Most modern processors use symmetric L1 caches
    \item \textbf{Computational Efficiency:} Reduces simulation count from 729 to 162, saving approximately 100+ hours
    \item \textbf{Focus:} Concentrates investigation on the actual bottleneck (L1D and L2)
\end{enumerate}

\subsubsection{Automation Infrastructure}

The sweep was executed using a parallel automation framework.
I wrote a small script to run all the simualations parallely.

\begin{itemize}
    \item Parallel Jobs: 20 simultaneous simulations
    \item Execution Time: Approximately 11 hours
    \item Logging: Real-time tracking in sweep\_progress.log
\end{itemize}

\subsection{Top 3 Optimal Configurations}

\subsubsection{Simple Merge Sort - Top 3 Configurations}

\begin{table}[H]
\centering
\begin{tabular}{clccc}
\toprule
\textbf{Rank} & \textbf{Configuration} & \textbf{IPC} & \textbf{L1D Miss} & \textbf{L2 Miss} \\
\midrule
1 & L1: 128KiB/4-way, L2: 1024KiB/4-way & \textbf{0.2895} & 0.227\% & 61.10\% \\
2 & L1: 128KiB/8-way, L2: 1024KiB/4-way & 0.2895 & 0.229\% & 60.51\% \\
3 & L1: 128KiB/16-way, L2: 1024KiB/4-way & 0.2895 & 0.230\% & 60.16\% \\
\bottomrule
\end{tabular}
\caption{Top 3 configurations for simple merge sort (ranked by IPC)}
\label{tab:top3_simple}
\end{table}

\textbf{Key Insight:} All top 3 configurations share maximum L1 size (128 KiB), maximum L2 size (1024 KiB), and minimum L2 associativity (4-way).

\subsubsection{Chunked Merge Sort - Top 3 Configurations}

\begin{table}[H]
\centering
\begin{tabular}{clccc}
\toprule
\textbf{Rank} & \textbf{Configuration} & \textbf{IPC} & \textbf{L1D Miss} & \textbf{L2 Miss} \\
\midrule
1 & L1: 128KiB/4-way, L2: 1024KiB/4-way & \textbf{0.2949} & 0.173\% & 49.19\% \\
2 & L1: 128KiB/8-way, L2: 1024KiB/4-way & 0.2949 & 0.173\% & 49.03\% \\
3 & L1: 128KiB/16-way, L2: 1024KiB/4-way & 0.2949 & 0.174\% & 48.86\% \\
\bottomrule
\end{tabular}
\caption{Top 3 configurations for chunked merge sort (ranked by IPC)}
\label{tab:top3_chunked}
\end{table}

\textbf{Key Insight:} Identical cache topology as simple sort, but achieves 1.9\% higher IPC, 24\% lower L1D miss rate, and 20\% lower L2 miss rate.

\subsection{Performance Analysis with Visualizations}

\subsubsection{Plot 1: L2 Miss Rate vs L2 Size}

\begin{figure}[H]
\centering
\includegraphics[width=0.85\textwidth]{plots/plot1_l2_miss_rate_vs_size.png}
\caption{L2 cache miss rate decreases significantly with larger cache sizes. The 1024KiB L2 achieves 48-61\% miss rate compared to 70-85\% for 256KiB, demonstrating L2 capacity is the dominant performance factor.}
\label{fig:plot1}
\end{figure}

\textbf{Analysis:} The working set (10MB) far exceeds even the largest L2 tested, but 1024KiB can hold enough of the merge operation's active data to provide meaningful benefits. L2 size shows the strongest correlation with performance among all parameters tested.

\subsubsection{Plot 2: IPC vs L1 Cache Size}

\begin{figure}[H]
\centering
\includegraphics[width=0.85\textwidth]{plots/plot2_ipc_vs_l1_size.png}
\caption{IPC shows moderate linear improvement with larger L1 caches. Chunked merge sort consistently outperforms simple merge sort across all L1 sizes, with the gap maintained at approximately 2\%.}
\label{fig:plot2}
\end{figure}

\textbf{Analysis:} L1 size has moderate impact on performance. Increasing from 32 KiB to 128 KiB yields approximately 1.5\% IPC improvement. The chunked algorithm maintains consistent advantage regardless of L1 configuration.

\subsubsection{Plot 3: L1D Hit Rate vs Associativity}

\begin{figure}[H]
\centering
\includegraphics[width=0.85\textwidth]{plots/plot3_l1d_hit_rate_vs_assoc.png}
\caption{L1D hit rates remain consistently high (greater than 99.7\%) regardless of associativity. This demonstrates that sequential merge sort access patterns do not cause significant conflict misses.}
\label{fig:plot3}
\end{figure}

\textbf{Analysis:} L1 associativity has minimal impact (less than 0.02\% IPC variation). Even 4-way associativity provides excellent hit rates due to sequential access patterns that avoid conflict misses.

\subsubsection{Plot 4: Simple vs Chunked Comparison}

\begin{figure}[H]
\centering
\includegraphics[width=0.85\textwidth]{plots/plot4_simple_vs_chunked_comparison.png}
\caption{Box plot comparison showing chunked merge sort consistently outperforms simple merge sort across all 81 configurations. The algorithmic advantage persists regardless of cache tuning.}
\label{fig:plot4}
\end{figure}

\textbf{Analysis:} Chunked shows 1.9-2.1\% IPC advantage across all configurations with consistently lower L1D and L2 miss rates throughout parameter space. Algorithm design trumps cache tuning for this workload.

\subsubsection{Plot 5: IPC Heatmap (L1 vs L2 Interaction)}

\begin{figure}[H]
\centering
\includegraphics[width=0.85\textwidth]{plots/plot5_ipc_heatmap.png}
\caption{Heatmap visualizing the interaction between L1 and L2 cache sizes. Vertical color gradients indicate L2 size dominates performance, while horizontal uniformity shows L1 size has secondary effect.}
\label{fig:plot5}
\end{figure}

\textbf{Analysis:} Strong vertical gradients indicate L2 size is primary factor. Weak horizontal gradients show L1 size is secondary. No significant interaction effects between L1 and L2 were observed.

\subsubsection{Plot 6: L2 Associativity Impact}

\begin{figure}[H]
\centering
\includegraphics[width=0.85\textwidth]{plots/plot6_l2_assoc_impact.png}
\caption{Surprising result: Higher L2 associativity slightly degrades performance. The 4-way configuration outperforms 16-way due to reduced tag comparison overhead with sequential streaming access patterns.}
\label{fig:plot6}
\end{figure}

\textbf{Analysis (Surprising Finding):} For 1024KiB L2 caches, 4-way associativity achieves the best performance. This counter-intuitive result is explained by: (1) Tag comparison overhead in 16-way requiring simultaneous comparison of 16 tags, adding complexity and latency, and (2) Sequential access optimization where simpler LRU on fewer sets (4-way) may evict more intelligently than complex policies on many ways.

\subsubsection{Plot 7: Performance vs Cache Budget}

\begin{figure}[H]
\centering
\includegraphics[width=0.85\textwidth]{plots/plot7_performance_vs_cache_budget.png}
\caption{IPC vs total cache size shows sublinear returns on cache investment. Quadrupling cache size yields only approximately 3\% performance improvement, indicating cache hierarchy saturation.}
\label{fig:plot7}
\end{figure}

\textbf{Cost-Benefit Analysis:} Minimum configuration (32KiB L1 + 256KiB L2 = 320 KiB total) achieves IPC of 0.286. Maximum configuration (128KiB L1 + 1024KiB L2 = 1280 KiB total) achieves IPC of 0.295, representing only 3.1\% improvement for 4x cache size increase. The cache hierarchy is saturated by this workload.

\subsubsection{Plot 8: Cache Miss Breakdown}

\begin{figure}[H]
\centering
\includegraphics[width=0.85\textwidth]{plots/plot8_cache_miss_breakdown.png}
\caption{Breakdown of cache misses by level. L2 misses dominate memory traffic, accounting for 48-85\% of L2 accesses. L1D misses are relatively low (less than 0.3\%), validating that L2 is the primary bottleneck.}
\label{fig:plot8}
\end{figure}

\subsubsection{Plot 9: Execution Time Comparison}

\begin{figure}[H]
\centering
\includegraphics[width=0.85\textwidth]{plots/plot9_execution_time_comparison.png}
\caption{Execution time across different configurations shows strong correlation with cache miss rates. Best configuration achieves 16.54s (simple) and 17.18s (chunked).}
\label{fig:plot9}
\end{figure}

\subsubsection{Plot 10: L1 Size × Associativity Interaction}

\begin{figure}[H]
\centering
\includegraphics[width=0.85\textwidth]{plots/plot10_l1_size_assoc_interaction.png}
\caption{Interaction plot confirming L1 associativity has negligible effect across all L1 sizes. Lines are nearly parallel, indicating no significant interaction between size and associativity.}
\label{fig:plot10}
\end{figure}

\subsubsection{Plot 11: Best vs Worst Configurations}

\begin{figure}[H]
\centering
\includegraphics[width=0.85\textwidth]{plots/plot11_best_vs_worst_configs.png}
\caption{Direct comparison of best (128KiB L1, 1024KiB L2, 4-way) vs worst (32KiB L1, 256KiB L2, 16-way) configurations. The 3\% performance range suggests cache hierarchy saturation.}
\label{fig:plot11}
\end{figure}

\subsection{Critical Insights}

\subsubsection{1. L2 Cache Size Dominates Performance}

L2 capacity is the single most important factor. A 1024KiB L2 provides 20-35\% lower miss rate than 256KiB, and this effect persists across all L1 configurations.

\subsubsection{2. L1 Associativity is Overrated for Sequential Access}

Associativity has minimal impact for sequential access patterns. The difference between 4-way and 16-way changes IPC by less than 0.02\%. Sequential merge operations avoid conflict misses, and lower associativity provides simpler, faster, cheaper hardware.

\subsubsection{3. Algorithm Design Matters More Than Hardware}

Chunked algorithm maintains 2\% advantage across all configurations due to better cache locality from streaming patterns and more predictable memory access for prefetchers. Software optimization complements hardware tuning.

\subsubsection{4. Cache Hierarchy Hits Fundamental Limits}

Only 3\% performance range across all configurations indicates the working set (10MB) far exceeds cache capacity. Even optimal caches cannot fully contain data, and memory bandwidth becomes the limiting factor.

\subsection{Design Recommendations}

\subsubsection{Optimal Configuration for Both Algorithms}

\textbf{Recommended:} L1: 128KiB/4-way, L2: 1024KiB/4-way

\textbf{Rationale:}
\begin{itemize}
    \item Maximizes L1 and L2 capacity (primary performance factors)
    \item Avoids unnecessary associativity overhead
    \item Best price/performance ratio: simpler hardware with equal performance
    \item Simple sort achieves IPC of 0.2895
    \item Chunked sort achieves IPC of 0.2949 (best observed)
\end{itemize}

\subsubsection{General Design Principles}

\begin{enumerate}
    \item Prioritize cache capacity over associativity for memory-intensive workloads
    \item Match L2 size to working set characteristics (at least 10\% of dataset for this workload)
    \item Consider algorithm-cache co-design rather than optimizing independently
    \item Test before over-engineering -- high associativity may not help and can hurt
\end{enumerate}

% ============================================================================
% SECTION 4: CONCLUSIONS
% ============================================================================
\newpage
\section{Conclusions}

\subsection{Complete Results Dataset}

For all 162 configurations tested (81 cache configurations $\times$ 2 algorithms), comprehensive performance data has been collected and organized in a structured CSV file located at \texttt{results/part2\_sweep/all\_results.csv}. This dataset contains 162 rows (one per simulation) and 20 columns capturing the following metrics for each configuration:

\textbf{Configuration Parameters (5 columns):}
\begin{itemize}
    \item sort\_type: Algorithm variant (simple or chunked)
    \item config\_name: Configuration identifier (e.g., L132KiB\_L1A4\_L2256KiB\_L2A4)
    \item l1\_size, l1\_assoc: L1 cache size and associativity
    \item l2\_size, l2\_assoc: L2 cache size and associativity
\end{itemize}

\textbf{Execution Metrics (3 columns):}
\begin{itemize}
    \item sim\_ticks: Total simulation cycles executed
    \item sim\_insts: Total instructions executed
    \item ipc: Instructions per cycle (primary performance metric)
\end{itemize}

\textbf{L1 Data Cache Metrics (4 columns):}
\begin{itemize}
    \item l1d\_accesses: Total L1D cache accesses
    \item l1d\_hits: L1D cache hits
    \item l1d\_misses: L1D cache misses
    \item l1d\_miss\_rate: L1D miss rate (fraction)
\end{itemize}

\textbf{L1 Instruction Cache Metrics (4 columns):}
\begin{itemize}
    \item l1i\_accesses: Total L1I cache accesses
    \item l1i\_hits: L1I cache hits
    \item l1i\_misses: L1I cache misses
    \item l1i\_miss\_rate: L1I miss rate (fraction)
\end{itemize}

\textbf{L2 Unified Cache Metrics (4 columns):}
\begin{itemize}
    \item l2\_accesses: Total L2 cache accesses
    \item l2\_hits: L2 cache hits
    \item l2\_misses: L2 cache misses
    \item l2\_miss\_rate: L2 miss rate (fraction)
\end{itemize}

This comprehensive dataset enables detailed analysis of cache behavior across the entire parameter space, facilitating identification of performance trends, optimal configurations, and parameter interactions. All 11 plots presented in Section 3 were generated from this consolidated results table. The complete CSV file is included in the submission package for independent verification and further analysis.

\subsection{Summary of Findings}

This comprehensive study of cache optimization for merge sort algorithms using gem5 reveals several critical insights about cache hierarchy design and algorithm-hardware interaction.

\subsubsection{Primary Findings}

\begin{enumerate}
    \item \textbf{Algorithm Design Trumps Cache Tuning:} The chunked merge sort consistently outperforms the simple version by 2\% across all 81 cache configurations, demonstrating that algorithmic optimization for cache locality is more impactful than hardware parameter tuning alone.

    \item \textbf{L2 Cache Capacity is Paramount:} L2 cache size is the dominant performance factor, with 1024KiB providing 20-35\% lower miss rates than 256KiB. For a 10MB working set, L2 should be at least 10\% of dataset size.

    \item \textbf{L1 Associativity Has Minimal Impact:} For sequential access patterns like merge sort, L1 associativity (4-way vs 16-way) changes performance by less than 0.02\%, validating that simpler, lower-associativity designs are sufficient.

    \item \textbf{Higher L2 Associativity Can Hurt:} Surprisingly, 4-way L2 associativity outperforms 16-way due to reduced tag comparison overhead with streaming access patterns.

    \item \textbf{Cache Hierarchy Saturation:} Only 3\% performance range across all configurations indicates the working set far exceeds cache capacity, limiting optimization potential through cache tuning alone.
\end{enumerate}

\subsubsection{Optimal Configuration Summary}

\textbf{Best Overall Configuration (Both Algorithms):}
\begin{itemize}
    \item L1: 128KiB, 4-way associative
    \item L2: 1024KiB, 4-way associative
    \item Simple IPC: 0.2895
    \item Chunked IPC: 0.2949 (+1.9\%)
\end{itemize}

\subsection{Design Implications}

\subsubsection{For Hardware Architects}
\begin{itemize}
    \item Prioritize cache capacity over associativity for memory-intensive workloads
    \item Consider that higher associativity may not improve performance and adds complexity
    \item Size L2 cache to match typical working set sizes of target applications
    \item Invest in prefetchers optimized for streaming access patterns
\end{itemize}

\subsubsection{For Software Developers}
\begin{itemize}
    \item Design algorithms with explicit consideration for cache hierarchy
    \item Use chunking and streaming patterns to improve locality
    \item Sequential access patterns are more cache-friendly than random access
    \item Blocking and tiling techniques can dramatically improve cache utilization
\end{itemize}

\subsection{Practical Recommendations}

For engineers working on memory-intensive applications:

\begin{enumerate}
    \item Start with algorithm optimization for cache locality before tuning hardware
    \item Profile first using cache simulators or hardware performance counters
    \item Prioritize capacity when in doubt -- choose larger cache size over higher associativity
    \item Test assumptions -- don't assume higher associativity always helps
    \item Balance performance gains against hardware complexity and cost
\end{itemize}

\subsection{Concluding Remarks}

This study demonstrates performance analysis using architectural simulation. The gem5 framework enables exploration of design spaces that giving insights that inform both hardware design and software optimization.

The key takeaway is that co-design of algorithms and cache hierarchies yields better results than optimizing either in isolation. The chunked merge sort's 2\% advantage persists across all cache configurations because it was designed with cache behavior in mind. This principle extends beyond sorting to any memory-intensive computational problem.

As memory access latencies continue to dominate processor performance, understanding and optimizing cache behavior becomes increasingly critical. This assignment provides a methodology for systematic cache analysis that can be applied to any computational workload.

% ============================================================================
% REFERENCES
% ============================================================================
\newpage
\begin{thebibliography}{9}

\bibitem{gem5docs}
gem5 Documentation. (2024).
\textit{Understanding gem5 statistics.}
Retrieved from https://www.gem5.org/

\end{thebibliography}

\end{document}
